\subsection{正式模型库与坐标}
正式研究包括九个目标、36 个预先指定的后验函数。为便于数值核对与解释，区分基础模型坐标 $r$、实际采样坐标 $q$ 和输出参数 $\theta$：$r=c+Lq$，$\theta=T_0(r)$；目标含 $\log|\det L|$ 及所需约束项。

\begin{table}[htbp]\centering\smalltable
\begin{tabularx}{\linewidth}{l r Y Y}\toprule
目标 & 维数 & 模型与数据 & 坐标及比较作用\\\midrule
G1 & 8 & 标准正态 & $c=0,L=I$；解析控制\\
G2 & 64 & 旋转高斯，原协方差条件数100 & 完整 Cholesky 白化；高维稠密计算\\
A1 & 64 & 潜在 AR(1)，$\rho=0.9$，单位观测噪声 & 后验均值平移与完整白化\\
L1 & 8 & logistic，$n=256$，先验标准差2.5 & 众数及观测信息坐标\\
L2 & 8 & logistic，$n=4096$，同一生成系数 & 似然成本及极稀有事件\\
H1 & 9 & 中心化漏斗，$v\sim N(0,9)$ & 原坐标，检验几何与路径敏感性\\
H2 & 9 & 非中心化漏斗 & 输出回原尺度，核随坐标改变\\
M1 & 8 & $\frac12N(-5e_1,I)+\frac12N(5e_1,I)$ & 双峰探索与分散初始化\\
W1 & 2 & 水井真实数据，3020条 & 截距、距离系数；二维参考\\\bottomrule
\end{tabularx}
\caption{当前 v4 的正式目标库。完整数据、矩阵、浮点配置及哈希随冻结协议保存。}
\end{table}

G2 的协方差为 $\Sigma=U\Lambda U^\mathsf T$，$\Lambda_{jj}=100^{(j-1)/63}$，$U$ 来自固定正态矩阵的 QR 分解。A1 的先验协方差为 $K_{ij}=0.9^{|i-j|}$，固定数据后 $\Sigma=(K^{-1}+I)^{-1}$、$m=\Sigma y$。二者在实际 $q$ 坐标均为标准正态，原尺度相关性不能直接解释本轮核的混合差异。

L1/L2 的设计含截距及七列标准化正态变量，系数先验为 $N(0,2.5^2)$，生成系数相同但数据不同。W1 为
$\Pr(y_i=1)=\sigma(\alpha+\beta d_i/100)$，在两个系数上用平坦先验。三种回归目标从零点开始，以 trust-exact 求众数；目标有限且解析梯度最大绝对值不超过 $10^{-7}$ 时接受，最多300步，$L$ 为观测信息逆矩阵的 Cholesky 因子。未做特征值修复，也未用参考样本构造坐标；该变换不把 Laplace 近似当作后验真值。

H1 为 $x_i\mid v\sim N(0,e^v)$，$i=1,\ldots,8$；H2 使用独立 $z_i\sim N(0,1)$ 并输出 $x_i=z_i e^{v/2}$。两者表示同一原尺度后验，但在各自坐标使用同一数值步长并不意味着同一原尺度转移核。

\subsection{独立开发选择与冻结配置}
步长在正式输入之外选择。RWM 候选为 $\{0.5,1,1.5,2.4,3.5,5\}/\sqrt d$；MALA 候选为 $\{0.1,0.3,0.6,1,1.6,2.4\}/d^{1/3}$。每候选两份四链输入，丢弃512步、保留1024步；仅在两次均有效时，以平均 ESJD/$d$ 最大者入选，平局取较小步长。该准则衡量局部移动，未优化 ESS 或稀有事件精度。

\begin{table}[htbp]\centering\smalltable
\begin{tabular}{lrrl}\toprule
目标 & MALA $h$ & RWM $s$ & 固定中心及线性变换\\\midrule
G1 & 0.8 & 0.848528137 & $0,I$\\
G2 & 0.4 & 0.3 & $0,\operatorname{chol}(\Sigma)$\\
A1 & 0.4 & 0.3 & $m,\operatorname{chol}(\Sigma)$\\
L1/L2 & 0.8 & 0.848528137 & 众数及观测信息\\
H1 & 0.480749857 & 1.66666667 & $0,I$\\
H2 & 0.769199771 & 0.8 & $0,I$\\
M1 & 0.8 & 0.848528137 & $0,I$\\
W1 & 1.26992084 & 1.69705627 & 众数及观测信息\\\bottomrule
\end{tabular}
\caption{冻结配置的显示精度。归档 JSON 提供完整精度；$h$ 对应噪声 $\sqrt{2h}$，$s$ 是 RWM 提议标准差。}
\end{table}

开发批次216项中212项完成，4项 H1 MALA 失败保留。H1 RWM 选择了候选上边界，说明配置仍受有限搜索范围约束。已记录开发任务内部成本226.917秒、坐标构造0.128秒；其中包含审计及写盘，未包含新进程导入，不能直接并入或扣除正式普通调用成本。

\subsection{正式网格、预算与资源}
每目标24份独立四链实际输入；同核的两个执行器及两设备读取同一输入。九种工作流为 CPU/CUDA 的四种 MH 组合，以及独立 CPU Pyro NUTS。两档预算每链保留1024或4096次转移，总计
\begin{equation}9\times9\times2\times24=3888\ \text{项主任务}.
\end{equation}
四链、时间执行器和设备重复不增加独立重复数。缓存探测限于 G1/G2/L2/W1，共256项、1024次调用；不计入正式后验拟合。

四链初值来自 $N(0,4I)$，M1 首坐标设为 $(-5,5,-5,5)$。MH 与 NUTS 共用初值，但 NUTS 使用独立随机流及预热。MH 丢弃512步、窗口32；quasi-DEER 每窗口最多2048轮，Picard最多 $N+512$ 轮。NUTS 用四个原生 spawn 子进程，每链预热1024步、目标接受率0.8、最大树深8，适应步长及对角质量矩阵，未使用 JIT。

正式 Windows 环境为原生 Windows 11、AMD Ryzen 7 9800X3D、RTX 5080，float64；Python 3.12.14、NumPy 2.2.6、SciPy 1.15.3、torch 2.13.0+cu130、Pyro 1.9.2、psutil 7.1.0，R 4.6.1、posterior 1.7.0。MH 进程 torch 线程4、interop线程1；NUTS每个子进程各为1。BLAS线程环境变量未统一设置，不能把这些数字解释为所有库的硬上限或独占物理核心预算。

存储/工作区估计上限2048 MiB、进程树 RSS 监督上限8 GiB、Windows Job提交内存上限12 GiB，三者计量不同。启动要求可用主存至少12 GiB、磁盘4 GiB，CUDA另需空闲显存3 GiB，运行磁盘底线1 GiB。资源不足暂停启动，数值失败、进程失败与未启动分别记录。没有总实验时限，也未通过重试或回退改变样本资格。

\subsection{函数、尺度与参考}
\begin{table}[htbp]\centering\smalltable
\begin{tabularx}{\linewidth}{l Y Y}\toprule
目标 & 原尺度上的函数集合 & 参考及其性质\\\midrule
G1/G2/A1 & 标准化首参数 $z$、$z^2$、$\ind(z>1)$ & 解析：$0,1,1-\Phi(1)$\\
L1/L2 & $\theta_1,\theta_2,\sigma(\theta_1),\ind(\theta_1>0)$ & 独立有限 NUTS；L2事件另作重要性参考\\
H1/H2 & $v/3,\tanh(v/3),\ind(x_1>0),\cos(x_1e^{-v/2})$ & 解析：$0,0,1/2,e^{-1/2}$\\
M1 & $\theta_1/\sqrt{26},\ind(\theta_1>0),\theta_2,\theta_2^2$ & 解析：$0,1/2,0,1$\\
W1 & $\alpha,\beta,\alpha^2,\beta^2,p(0),p(100),\ind(\beta>0)$ & 冻结求积点；v4另给确定数值包络\\\bottomrule
\end{tabularx}
\caption{全部36函数。各尺度 $s_f=1$，因此阈值表示函数原单位上的绝对误差，而非自动统一的相对精度。}
\end{table}

对方法 $a$ 和预算 $N$，在数值有效且函数有限的重复集合 $V$ 上计算
\begin{equation}\label{eq:loss}
\widehat E_{a,f}(N)=\frac{1}{|V|}\sum_{r\in V}
\frac{(\widehat\mu_{a,f}^{(r)}(N)-\mu_f^{\rm ref})^2}{s_f^2}.
\end{equation}
解析参考时可解释为有效输出条件下的 MSE 估计；有限参考时称为参考平方差。比较两方法固定共同有效集合，同时报告全部24份的失败分母，不以 $\hat R$ 达标筛掉困难拟合。预算点上的最大函数平方差 $\leq0.01$ 对应 RMSE $\leq0.1$，只对指定函数与尺度有意义。误差估计的 MCSE 判据与这一误差阈值分开，不把短链 ESS 当作真实 MSE。

主分析使用9999次逐点 BCa bootstrap，单位为整份四链独立输入，至少20份共同有效重复才计算区间；退化情况保留未定，不作全体比较的同时置信声明。秩正态化与折叠 $\hat R$、bulk/tail ESS 由 posterior 计算\cite{vehtari,posterior}；常量函数不填成 $\hat R=1$。预算仅两个点，报告已测预算的误差与耗时，不插值为精确达标时间。

L1/L2 每模型四份独立四链 NUTS 参考，每链预热1024、保留16384步，共262144个相关状态；MCSE取批均值与拟合间估计的较大者。无记录发散不能排除共同偏差。L2原符号事件全零，其不确定性保持未定。W1 原参考以12个保存的浮点求积配置检查稳定性；这些变化不是认证误差界。v4的 L2 重要性参考和 W1 确定包络单独报告，不覆盖原参考身份。

\subsection{成本边界与比较单位}
\begin{table}[htbp]\centering\smalltable
\begin{tabularx}{\linewidth}{l Y Y}\toprule
口径 & 包含内容 & 回答的问题\\\midrule
缓存执行 & 已准备内核的执行，显式设备同步 & 同一指定计算能否更快完成？\\
普通批量工作流 & 进程启动、输入/目标准备、预热/丢弃、采样、变换、正常结果和 R 诊断 & 获得可分析输出的实际批量成本\\
研究审计 & 再含独立路径参考、保护、完整证据归档 & 取得可核验研究材料的成本\\
R 前端技术成本 & 最外层 R 启动/依赖、结果传回与转换的实测 & R 使用者首次与重复调用的边界\\\bottomrule
\end{tabularx}
\caption{成本口径相互嵌套或属于另测边界，不能直接相加。开发调参、整库环境安装和全局归档上传另列。}
\end{table}

同设备、同核、同输入配对的速度比为 $S=T_{\rm seq}/T_{\rm par}$，$S>1$ 表示时间执行更快。不同核需结合函数误差与诊断解释。跨 Mac/JAX 和 Windows/torch 的差异同时包含硬件、后端及实现，不归因于算法单一因素。批处理向量化不等于已证明物理核心有效利用；CPU/CUDA计时包含实际同步及主机控制成本。
